\begin{abstract}
Using a digital wallet in Colombia does not mean a woman can complete a transfer across borders or control how the money is used. We study migrant women's reported handling of remittances and participation in domestic digital payments, alongside existing interviews, focus groups, expert interviews and a five-country institutional comparison. Migrant accounts describe routine domestic use alongside reliance on intermediaries for some cross-border journeys, as well as variation by corridor. The study distinguishes reported use from the ability to open, operate and protect an account. It does not estimate the effect of Colombia's recent interoperable payments rollout.
\par\smallskip
\emph{Working-draft note: the quantitative findings sentence and final abstract await accepted exhibits.}
\end{abstract}
\noindent\textbf{Keywords:} digital payments; remittances; migrant women; financial inclusion; digital public infrastructure; Colombia.
